\section*{第 \ref{ch:b3:schrodinger-three-dimensions} 章——三维薛定谔方程}

\begin{solution}{exo:b3:schrodinger-three-dimensions:1}
(a) $\int|\psi|^2\dd^3r = |A|^2\,4\pi\int_0^\infty r^2\eu^{-r^2/2
\sigma^2}\dd r = |A|^2(2\pi\sigma^2)^{3/2}$: $A = (2\pi\sigma^2)^{
-3/4}$。(b) 由各向同性，$\langle r^2\rangle = 3\langle x^2\rangle =
3\sigma^2$，故 $\Delta x = \sigma$。(c) $\vect\jmath = \vect 0$：
波函数为实函数。(d) $\langle\vect p\rangle = \hbar\vect k_0$
且 $\vect\jmath = \rho\,\hbar\vect k_0/m$：同一团云，此刻
正在漂移。
\end{solution}

\begin{solution}{exo:b3:schrodinger-three-dimensions:2}
$3\,(1)$, $6\,(3)$, $9\,(3)$, $11\,(3)$, $12\,(1)$, $14\,(6)$ ---
此后直到 $17$ 之前皆无。能级 $27$：$(3,3,3)$ 以及
$(5,1,1)$ 的置换——共四个态，因为 $27 = 9+9+9 = 25+1+1$。立方
对称性仅置换坐标轴，且无置换能将 $(3,3,3)$ 与
$(5,1,1)$ 相联：这一巧合是算术性的（
同一数作为三个平方数之和有两种表示），而非几何性的。
\end{solution}

\begin{solution}{exo:b3:schrodinger-three-dimensions:3}
(a) $\vect\jmath = |A|^2\hbar k/m\,\vect e_x$。(b) 为零：波函数为实
函数。(c) $j = (|A|^2\hbar k/m)(1 - r^2)$：入射通量减去
反射通量——交叉项相消。(d) $j = \rho\,\hbar
m/Mb$，其中 $\rho = 1/2\pi$（每单位角度）：一个永久
环流，即持久电流的微观前身。
\end{solution}

\begin{solution}{exo:b3:schrodinger-three-dimensions:4}
系数为 $\pi^2(\hbar c)^2/2\mu c^2 = \qty{3.76}{eV.nm^2}$。
(a) 当 $R = \qty{3}{nm}$ 时：\qty{0.42}{eV}。(b) $E = \qty{2.16}{eV}$：
$\lambda = 1240/2.16 = \qty{575}{nm}$，黄绿色。(c) $R$ 越小，
$E$ 越大：向蓝端偏移。(d) 在块体中，限域项
消失：带隙固定为 \qty{1.74}{eV}，$\lambda = \qty{713}{nm}$，
无论样品如何，都是同样的深红色。
\end{solution}

\begin{solution}{exo:b3:schrodinger-three-dimensions:5}
(a) 单值性：$\psi(\varphi + 2\pi) = \psi(\varphi)$ 要求
$\eu^{2\pi\iu m} = 1$。(b) $E_m = \hbar^2m^2/2Mb^2$；态 $\pm
m$ 反向环流，并经镜面反射（或
时间反演）互换——环的一种对称性，因而导致简并。(c)
电流环 $I = q\hbar m/2\pi Mb^2$ 携带磁矩
$\mu = I\pi b^2 = m\,q\hbar/2M$——以 $q\hbar/2M$ 为步长量子化，
即\cref{ch:b3:quantum-angular-momentum}中的玻尔磁子模式。
(d) $\hbar^2/2m_{\text{e}}b^2 = \qty{1.9}{eV}$；六个电子
占据 $m = 0, \pm1$，第一激发 $\pm1 \to \pm2$ 需要
$(4-1) \times 1.9 \approx \qty{5.8}{eV}$，即 $\lambda \approx
\qty{210}{nm}$——一个环所给出的恰是紫外
附近区域，此外别无其他。
\end{solution}

\begin{solution}{exo:b3:schrodinger-three-dimensions:6}
(a) 取 $n_x = 0..n$ 及 $n_y = 0..n - n_x$：$\sum_{n_x}(n - n_x
+ 1) = (n+1)(n+2)/2$。(b) 计入自旋：$2, 6, 12, 20$；累计 $2,
8, 20, 40$。(c) $2$、$8$、$20$ 为幻数；超出这些数之后的失效说明
核势阱并非谐振子型——底部更平坦，而最重要的是
并具有很强的自旋–轨道耦合，它使
壳层重排为 $28, 50, 82$。(d) 壳层结构需要\emph{成簇的}
能级，它们由间隙隔开；而盒子的能级分布不规则，
没有大间隙，因此不会出现类似原子核的稳定性规律。
\end{solution}

\begin{solution}{exo:b3:schrodinger-three-dimensions:7}
(a) $u_n = \sqrt{2/R}\sin(n\pi r/R)$, $E_n = n^2\pi^2\hbar^2/2MR^2$.
(b) $E_1 = \pi^2(\hbar c)^2/2Mc^2R^2 = 384210/(2 \times 939 \times
25) = \qty{8.2}{MeV}$。(c) 量级正确：原子核中的核子处于
间隔为 MeV 量级的量子态，正如其能谱所示。(d) $|u_1|^2$
峰值位于 $r = R/2$：径向概率在半途最大
（体积元中的 $r^2$ 隐含在 $u = r\varphi$ 中，将
极大值推离中心），而一维盒子的 $|\varphi|^2$ 峰值位于
中心。
\end{solution}

\begin{solution}{exo:b3:schrodinger-three-dimensions:8}
(a) 令 $N = 2N_{\text{态数}}(E_{\text{F}})$，并将计数反解。
(b) 铜：$E_{\text{F}} = \qty{7.1}{eV}$。(c) $v_{\text{F}} =
\sqrt{2E_{\text{F}}/m_{\text{e}}} = \qty{1.6e6}{m/s}$;
$T_{\text{F}} = E_{\text{F}}/k_{\text{B}} \approx \qty{8e4}{K}$. (d)
电子是全同费米子：泡利不相容原理
（见\cref{ch:b3:identical-particles}）至多允许每个轨道
态容纳两个电子，因此电子群必须堆叠到电子伏特量级的能量。
\end{solution}

\begin{solution}{exo:b3:schrodinger-three-dimensions:9}
(a) $k = \sqrt{2M(V_0 - |E|)}/\hbar$, $\kappa = \sqrt{2M|E|}/\hbar$;
由 $u'/u$ 在 $R$ 处的连续性，得 $k\cot kR = -\kappa$。(b) 当 $|E|
\to 0$，$\kappa \to 0$：$\cot kR = 0$，$kR = \pi/2$，于是 $V_0 =
\hbar^2k^2/2M = \pi^2\hbar^2/8MR^2$。(c) 取 $\mu c^2 =
\qty{469}{MeV}$, $R = \qty{2.1}{fm}$: $V_{0,\min} = \qty{23}{MeV}$.
(d) 硬壁条件 $u(0) = 0$ 使三维 $s$ 态成为一维
\emph{奇}态的类比，它必须在阱内容纳四分之一波长：
浅阱做不到，而在一维情形下，无节的偶态总能
形成束缚态。
\end{solution}

\begin{solution}{exo:b3:schrodinger-three-dimensions:10}
(a) $\kappa = \sqrt{2\mu c^2B}/\hbar c = \sqrt{2 \times 469 \times
2.22}/197.3 = \qty{0.231}{fm^{-1}}$：尾部长度 $1/\kappa =
\qty{4.3}{fm}$，为力的作用范围的两倍。(b) $k =
\sqrt{2\mu(V_0 - B)}/\hbar$ 及 $\cot kR = -\kappa/k$。(c) 迭代：
$kR = 1.83$, $k = \qty{0.87}{fm^{-1}}$, $V_0 - B = (\hbar ck)^2/\mu
c^2\cdot\tfrac12 \approx \qty{31}{MeV}$: $V_0 \approx \qty{34}{MeV}$.
(d) 内部：$\int_0^R\sin^2kr\,\dd r \approx \qty{1.19}{fm}$；
外部：$\sin^2(kR)/2\kappa \approx \qty{2.0}{fm}$：约 $63\%$ 的
概率位于\emph{阱外}——氘核大多
栖居于其束缚力已经消失的区域，这是一种
纯粹的量子晕。
\end{solution}

\begin{solution}{exo:b3:schrodinger-three-dimensions:11}
(a) $\dd\langle x\rangle/\dd t = \int x\,\partial_t\rho\,\dd^3r =
-\int x\operatorname{div}\vect\jmath\,\dd^3r = \int j_x\,\dd^3r =
\langle p_x\rangle/m$（分部积分，边界项为零）。(b)
对 $\langle\vect p\rangle = -\iu\hbar\int\psi^*\vect\nabla
\psi$ 求导，两次利用方程；动能项经分部积分相消，
留下 $-\int|\psi|^2\vect\nabla V$。(c) 若 $V$ 至多为二次式：
则 $\langle\vect\nabla V(\vect r)\rangle = \vect\nabla V(\langle
\vect r\rangle)$ 严格成立。(d) 经双缝后，波包分为两
瓣；$\langle\vect r\rangle$ 沿对称轴下滑，而
粒子几乎从不落在那里：平均值描述的是
系综的质心，而非任何一个粒子的轨迹。
\end{solution}

\begin{solution}{exo:b3:schrodinger-three-dimensions:12}
(a) 动能：尺度为 $1/a$ 的梯度给出 $\hbar^2/2m_{\text{e}}
a^2$；势能：电子云位于距离 $\sim a$ 处，给出
$-ke^2/a$（系数恰好精确为 $1$）。(b) 由 $\dd E/\dd
a = 0$ 得 $a = \hbar^2/m_{\text{e}}ke^2 = a_0$；$E(a_0) =
-ke^2/2a_0 = \qty{-13.6}{eV}$。(c) 在 $a_0$ 以下时，动能代价
按 $1/a^2$ 增长，快于 $-1/a$ 的收益：局域化
需付出不确定性原理所要求的代价，原子漂浮于
收支平衡的尺度上。(d) 将 $ke^2$ 换为 $Gm_{\text{e}}m_{\text{p}} =
\qty{1.0e-67}{J.m}$：$a_{\text{引力}} = \hbar^2/m_{\text{e}}G
m_{\text{e}}m_{\text{p}} \approx \qty{1.2e29}{m}$——比
可观测宇宙还大：引力太弱，不足以闭合
量子轨道，它造出的是恒星，而非原子。
\end{solution}

\begin{solution}{pb:b3:schrodinger-three-dimensions:1}
\textbf{1.} $s$ 波：$u = rR$ 在球内服从
一维自由方程，$u(0) = u(R) = 0$ 时 $u \propto \sin(\pi r/R)$，能量
$\pi^2\hbar^2/2\mu R^2$。
\textbf{2.} 盒中基态能量为正，且随
盒子变小而增大——局域化总要付出动能代价
（不确定性原理）；它只能使能隙增大。
\textbf{3.} 取 $\pi^2(\hbar c)^2/2\mu c^2 = \qty{3.76}{eV.nm^2}$：
$0.04$、$0.24$、$0.94$、\qty{3.8}{eV}——在 $R \approx \qty{2}{nm}$
时，修正量几乎与能隙本身相当。
\textbf{4.} 吸引作用降低电子–空穴对的能量：发射
光红移。它按 $1/R$ 变化，与限域的 $1/R^2$ 相比：
对小量子点，盒子项占主导，忽略该项的近似反而更好。
\textbf{5.} $1240/1.74 = \qty{713}{nm}$：处于人眼可见的红端。
\textbf{6.} 单一颜色来自 $10^{15}$ 个发光体，证明合成工艺使它们
全部具有\emph{相同尺寸}——单分散性，
这一化学成就是其物理背后的根基。
\textbf{7.} $1.97$、$2.34$、\qty{2.70}{eV}。
\textbf{8.} 限域能 $0.23$、$0.60$、\qty{0.96}{eV}：$R =
\sqrt{3.76/\Delta E}$ 给出 $4.0$、$2.5$、\qty{2.0}{nm}。
\textbf{9.} 直径 $8.1$、$5.0$、\qty{4.0}{nm}：约跨越 $13$、
$8$ 和 $7$ 个晶格间距。
\textbf{10.} 蓝色量子点最小，此处 $E$ 对 $R$ 的依赖最为
陡峭：相差一个晶面的误差使颜色偏移最大——
小量子点最不容许误差。
\textbf{11.} 求导：$\dd E/\dd R = -2\Delta E_{\text{约束}}/R
= -\pi^2\hbar^2/\mu R^3$；在 $R = \qty{2.5}{nm}$ 处：$2 \times 0.60/2.5
= \qty{0.48}{eV/nm} = \qty{480}{meV/nm}$。
\textbf{12.} $\pm 5\% = \pm\qty{0.125}{nm}$：$\Delta E = \pm
\qty{60}{meV}$，即 $\Delta\lambda = \lambda^2\Delta E/hc \approx
\pm\qty{14}{nm}$——展宽约 $\sim\qty{28}{nm}$，恰好处于
显示器的容差：百分之五恰是\emph{边界}
，即可接受化学的范围。
\textbf{13.} 因为颜色纯度要求在
单晶面层级上控制尺寸，覆盖 $10^{15}$ 个粒子——这正是
叶基莫夫、布鲁斯与巴文迪实现的可控生长，
并获 2023 年诺贝尔奖委员会引用。
\textbf{14.} 量子点只能在其自身（最低）能隙处发射；吸收
需要至少该能量，因此泵浦光须更蓝。多余的能量
以晶格振动的形式弛豫，变为热量。
\textbf{15.} $1 - 460/630 = 27\%$ 的泵浦光子能量加热
屏幕。
\textbf{16.} 发光能级以上，量子点的能谱很密，
（盒子态众多）：吸收在宽频带上都能成功；随后
激发迅速跌落至最低激发态，发射来自这一单一能级：
谱线很窄。
\textbf{17.} 所有量子点都能欣然吸收同一紫外线，
（宽频带），而每种尺寸各发其色：一盏灯，
多种标记——与染料化学恰好相反。
\textbf{18.} $\tfrac43\pi(4.0)^3/(0.3)^3 \approx \num{e4}$ 个原子：
一万个原子共享一套离散的类氢能级，
能级——“人造原子”实至名归。
\textbf{19.} 晶体边界：半导体晶粒的限域壁
取代了原子核的吸引，成为束缚并使电子量子化的
因素。
\textbf{20.} $\pi^2(\hbar c)^2/2Mc^2R^2 = 384210/(2 \times 939
\times 9) \approx \qty{23}{MeV}$：核物理使用 MeV，
正是因为飞米尺度的盒子如此。
\textbf{21.} 对电子而言，盒子估算须作相对论处理；
老实说，$E \sim \pi\hbar c/R \approx \qty{200}{MeV}$——然而 β 衰变
中的电子只以几 MeV 射出：电子不能是原子核的组成部分，
这一论证埋葬了旧的核–电子模型，
并为中微子的提出铺平了道路。
\textbf{22.} $\pi^2\hbar^2/2ML^2 \approx \qty{2e-70}{J}$：
比热噪声低三十九个数量级——人并非量子
限域的。
\textbf{23.} $\pi^2\hbar^2/2m_{\text{e}}R^2 = m_{\text{e}}c^2$ 时，
$R = \pi\hbar/\sqrt2\,m_{\text{e}}c \approx \qty{0.9}{pm}$，即
康普顿尺度：此处的限域能足以产生电子–正电子对，
单粒子方程于是让位于
量子场论。
\textbf{24.} $E_{\text{约束}} \propto 1/MR^2$：较轻的粒子与
较小的盒子更具量子性，一个公式便可说明。
\textbf{25.} $E_{\text{g}} + \pi^2\hbar^2/2\mu R^2$ 将 $4.0 \to$
红色，$2.5 \to$ 绿色，$2.0\,\unit{nm} \to$ 蓝色：
盒中粒子经化学家精确调至纳米尺度，点亮了
店铺橱窗——限域量子化已化身为消费电子产品。
\end{solution}
